% !TeX root = ../main.tex
\section{Generation 1: the published result, re-analysed and re-graded}
\label{sec:gen1}

Before asking why look-ahead stopped working, we establish that it worked --- and that it
worked for the reason claimed, not through an artifact of the analysis or of the grader.

\subsection{The published comparison}
The ICLR paper compared the best model of each look-ahead arm: $\K{=}5$ iteration~7 against
$\K{=}0$ iteration~4, unpaired. On its own scores that comparison gives $+0.206$ on the
composite (Welch $p = 0.023$), $+0.191$ on Q2 ($p = 0.011$) and $+0.221$ on Q1
($p = 0.063$). The paper reported this accurately, including the observation that only Q2
reached significance. Two features of the design nonetheless weaken it: selecting each arm's
maximum before comparing them selects on the outcome, and an unpaired test throws away the
persona structure that dominates the variance.

\subsection{Re-analysed: paired, at matched iterations}
Replacing that with the test of \S\ref{sec:test} --- persona-paired, at matched iteration
counts, across all seven iterations --- makes the effect \emph{stronger}. On its original
GPT-3.5 scores, $\K{=}5$ leads at \textbf{all seven} matched iterations on the \QQ{} composite,
with mean $\dlt = -0.206$; all $21$ (iteration $\times$ metric) cells favour look-ahead, and
$7$ survive Holm correction. The three significant composite cells are iteration~3
($\dlt = -0.185$, $\dz = -0.23$), iteration~5 ($-0.332$, $\dz = -0.32$) and iteration~7
($-0.418$, $\dz = -0.47$).

This matters for how the rest of the chapter is read. The original finding is not a fragile
artifact that a better test dissolves; a better test endorses it more strongly than the
authors claimed. Whatever explains generation~3's disagreement, it is not that generation~1
used a weak statistic.

\subsection{Re-graded: is it the judge?}
\label{sec:regrade}
The remaining mundane explanation is the grader. GPT-3.5 filling in a regex-parsed rubric is
a weaker and noisier instrument than \primaryjudge{} filling in a JSON schema, and a weaker
instrument could plausibly reward the more polished-looking openings that look-ahead selects
for, without those openings being better therapy.

We test this directly by re-scoring generation~1's conversations with the generation-3 oracle
(\S\ref{sec:one-grader}). Table~\ref{tab:gen1-graders} shows the result: the effect survives.

\begin{table}[t]
\centering\small
\begin{tabular}{@{}lrrrr@{}}
\toprule
\textbf{Iter.} & \multicolumn{2}{c}{\textbf{GPT-3.5} (original)} & \multicolumn{2}{c}{\textbf{\primaryjudge{}} (re-graded)} \\
\cmidrule(lr){2-3}\cmidrule(lr){4-5}
 & $\dlt$ & $\dz$ & $\dlt$ & $\dz$ \\
\midrule
1 & $-0.053$ & $-0.05$ & $-0.038$ & $-0.05$ \\
2 & $-0.201$ & $-0.21$ & $-0.138$ & $-0.18$ \\
3 & $-0.185^{*}$ & $-0.23$ & $-0.205^{*}$ & $-0.29$ \\
4 & $-0.017$ & $-0.02$ & $-0.074$ & $-0.11$ \\
5 & $-0.332^{*}$ & $-0.32$ & $-0.090$ & $-0.13$ \\
6 & $-0.240$ & $-0.24$ & $-0.094$ & $-0.13$ \\
7 & $-0.418^{*}$ & $-0.47$ & $-0.284^{*}$ & $-0.37$ \\
\midrule
mean & $\mathbf{-0.206}$ & $-0.21$ & $\mathbf{-0.132}$ & $-0.18$ \\
$\K{=}5$ ahead & \multicolumn{2}{c}{$7/7$} & \multicolumn{2}{c}{$7/7$} \\
\bottomrule
\end{tabular}
\caption{Generation~1 under both graders, on the \emph{same} $1{,}344$ conversations.
$\dlt = \K_0 - \K_5$ on \QQ{}, persona-paired, $n = 96$ per cell; negative favours look-ahead.
$^{*}$ Holm-corrected $p < 0.05$. The effect attenuates by about a third under the stronger
grader but does not reverse or vanish.}
\label{tab:gen1-graders}
\end{table}

Under \primaryjudge{}, $\K{=}5$ still leads at all seven matched iterations, mean
$\dlt = -0.132$, with $20$ of $21$ metric cells favouring look-ahead and $6$ surviving Holm
correction. The two graders agree on the same conversations at a mean Pearson
$r = 0.774$ (range $0.542$--$0.867$), with \primaryjudge{} scoring $+0.269$ higher on average
--- a level offset, not a reordering.

The attenuation from $-0.206$ to $-0.132$ is worth noting rather than dismissing: roughly a
third of the published effect size does appear to have been grader-specific. But two-thirds of
it is not, and the direction and consistency are unchanged. \textbf{The grader change cannot
explain a reversal, because on generation~1's own conversations the modern grader reproduces
the original conclusion.}

\subsection{The detail that turns out to matter}
One feature of generation~1 is invisible in the $\K$ contrast but becomes central in
\S\ref{sec:discussion}: under the common grader, \textbf{generation~1's $\K{=}0$ arm barely
learns at all}. Its untrained base scores $3.865$ on \QQ{}; after seven PTO iterations the
myopic arm reaches $3.887$, a gain of $+0.023$ --- well inside grader noise --- with a fitted
slope of $+0.011$ per iteration. The $\K{=}5$ arm over the same seven iterations climbs to
$4.171$ at $+0.030$ per iteration.

So in the generation that made look-ahead's reputation, the comparison was not
``good training versus better training''. It was ``training that did not work versus training
that did''. Look-ahead was not adding a margin on top of a functioning myopic signal; it was
supplying a signal where the myopic one had failed. Figure~\ref{fig:trajectories} (left panel)
shows this directly --- the $\K{=}0$ trajectory tracks the untrained baseline for all seven
iterations.
